\section{Introduction}\label{sec:introduction}

How many integer variables are intrinsically necessary to approximate a
nonlinear graph to a prescribed accuracy? A formulation can use many
interval selectors, a few integer variables with large ranges, continuous
lifting variables, or long coefficients. An efficient encoding of one
chosen partition gives an upper bound, but it does not establish what is
possible with all convex mixed-integer formulations. We study this minimum
and identify geometric and algebraic quantities that govern it.

For a continuous map $F:D\to\R^m$ and a prescribed error body $K$, an
admissible projected relaxation $R$ must satisfy
\[
 \{(x,F(x)):x\in D\}\subset R
 \subset\{(x,w):x\in D,\ w-F(x)\in K\}.
\]
Thus the approximation retains every input-output pair and controls the
output error at the same input. We minimize the number of general integer
coordinates in an arbitrary convex lift of any such $R$, denoting the
result by $\pconv$. Integer ranges, finite continuous dimension, and real
coefficients are unrestricted; the lift need not be closed. The comparison
quantity $\pbin$ restricts the lift to linear constraints and binary
coordinates, so $\pconv\le\pbin$. When we give a polynomial-time algorithm
producing a formulation of polynomial rational encoding length, we report
its binary count separately as $\pout$.

The minimum dimension of a convex integer representation is already the
\emph{MICP rank} of \citet[Definition 4.3]{lubin2022}, who also prove a
parity-based lower bound. Our accuracy-dependent question minimizes over
all intermediate graph relaxations $R$, and allows the broader class of
nonclosed convex lifts. On the constructive side, logarithmic encodings
of bounded polyhedral unions and compact quadratic approximations are established
\citep{vielma2018,beach2022}. The issue addressed here is whether such
constructions use close to the fewest integer coordinates allowed by the
geometry, even when compared with general integers and arbitrary convex
constraints.

This comparison separates three questions that matter in formulation
design: the integer count required by accuracy, the benefit of using
nonlinear convex constraints or unbounded integer ranges, and the cost of
writing down a formulation. In particular, an asymptotically optimal count
can coexist with an impractically long description. Our results provide
benchmarks for these resources and identify when one discretization can
serve several nonlinear expressions. They concern approximation of the
entire graph on the stated domain; restricting that domain by additional
nonlinear equations can change the lower bounds. Small integer dimension
also does not imply fast optimization or an ideal linear relaxation.

\subsection{Main contributions}

The principal developments are an exact quadratic precision law, a finite
quadratic comparison uniform over coefficients and outputs, and transfers
from scalar approximation to minimum integer dimension.

\paragraph{An algebraic precision law for quadratic systems.}
Let $F$ be a fixed real quadratic system on a full-dimensional bounded
box $B\subset\R^n$, and let $\mathcal H$ be the span of its Hessians. With componentwise
absolute error $\eps>0$, \cref{thm:ncrank} proves
\begin{equation}\label{eq:intro-ncrank}
 \pconv(\eps),\ \pbin(\eps)
 =\tfrac12\ncrank(\mathcal H)\log_2(1/\eps)+O_{F,B}(1)
 \qquad(\eps\downarrow0).
\end{equation}
Both minima have this expansion; they need not agree at a fixed tolerance.
Here noncommutative rank is an established matrix-space invariant
\citep{ggow2020,iqs2018}. One useful characterization is
\[
 \ncrank(\mathcal H)
 =n-\max_{U\subseteq\R^n}
       \bigl(\dim U-\dim\mathcal H(U)\bigr),
 \qquad
 \mathcal H(U)=\operatorname{span}\{Hu:H\in\mathcal H,\ u\in U\},
\]
where the maximum is over linear subspaces. The maximum records the largest loss
of dimension forced by all Hessians on a common subspace; noncommutative
rank subtracts that loss from the input dimension. For symmetric
Hessian spaces the real form of this characterization follows from the
usual complex one, as shown in \cref{lem:symmetric-shrink}.

The joint invariant is essential. For the map $(u,v)\mapsto u\times v$
on six inputs, every scalar combination of the three outputs has Hessian
rank at most four, whereas the Hessian space has noncommutative rank six.
The coefficient in \eqref{eq:intro-ncrank} is therefore three, which no
single scalar combination detects. To the best of our knowledge, the
identification of half the noncommutative rank with the minimum
integer-dimension precision coefficient for quadratic graph approximation
is new. The algebraic invariant and its algorithms are not new; the
contribution is the lower bound for all admissible convex lifts together
with a matching binary linear construction.

The proof makes the connection explicit. Contacts with the exact graph
whose integer witnesses have the same parity have admissible midpoints.
Bounds on the covariance and volume of these contacts limit how much of
the input domain one parity class can cover. Symmetric shrinking supplies
a coordinate decomposition with discretization depths proportional to
$\log_2(1/\eps)$, with coefficients in $\{0,1/2,1\}$ summing to
$\ncrank(\mathcal H)/2$. Shared product relaxations then attain the lower
coefficient. Scalar quadratic rank, one-sided inertia, and fractional
vertex-cover formulas for bilinear interaction graphs are related exact
cases. Local/global Hessian-rank bounds, constant scalar Hessian rank,
and positive perspectives extend parts of the analysis beyond quadratics
(\cref{thm:smooth-ranks,thm:constant-rank,prop:perspective}).

\paragraph{A finite-accuracy benchmark and a rational construction.}
An asymptotic coefficient does not describe accuracy thresholds caused by
small or highly unequal coefficients. For quadratic outputs with Hessians
$H_j$ on $[0,1]^n$ and positive component tolerances $\eps_j$, define
\[
 D_* = \max\{\det P:0\preceq P\preceq I,
       \ \operatorname{tr}(H_jPH_jP)\le\eps_j^2\ \forall j\},
 \qquad \Phi=-\tfrac12\log_2 D_*.
\]
\Cref{thm:finite-covariance} places both minima within
$O(n\log(n+1))$ of $\Phi$, with constants independent of the coefficients,
tolerances, and number of outputs. \Cref{thm:rational-finite} constructs,
for rational data, a rational MILP with
$\pout\le\pconv+O(n\log(n+1))$ in polynomial time and encoding length.
The determinant constraints need not be Euclidean convex; the algorithm
uses their geodesic convexity, rational approximations, and explicit
feasibility repairs.

The contribution is the common finite benchmark for arbitrary convex
integer lifts and binary linear constructions, including the rational
algorithm. Anisotropic interpolation, determinant optimization, and
geodesic optimization have their own established theories, discussed
below. Further results give finite certificates, output-body and
nonlinear-input reductions, and sharper positive semidefinite block and
diagonal bounds. A thin-domain example shows why a containing box can
misrepresent a correlated input domain. A complexity reduction also limits
polynomial-time additive approximation of the minimum count: achieving
error $O(n^{1-\delta})$ for any fixed $\delta\in(0,1)$ on all rational
quadratic instances would imply $P=NP$.

\paragraph{Scalar and vector comparisons with controlled descriptions.}
For every continuous convex scalar function on a compact interval and
every positive absolute tolerance, \cref{lem:scalar-chords} proves
\[
 \pconv\le\pbin\le\pconv+2.
\]
Thus general integer ranges and nonlinear convex constraints save at most
two integer coordinates in this setting. The bound uses real coefficients
and does not control continuous formulation size. A separate result,
\cref{thm:convex-hybrid}, constructs a rational MILP within eleven binaries
of $\pconv$ for every rational convex polynomial on $[0,1]$ given in dense
encoding and every positive rational tolerance. Its running time and total
encoding length are polynomial in the input length, including the tolerance.
This requires
compact access to the selected intervals even when their number is
exponential in the input length. An accuracy-dependent truncated curvature
measure and a hybrid of greedy segmentation and curvature integration
provide that access.

For several convex outputs of one input, shared breakpoints are already
available \citep{lyu2026}. Let $r$ be the dimension of the outputs' span
modulo affine functions, called their \emph{curvature rank}. For $r\ge1$,
our comparison bounds the additional binary count by $\lceil\log_2r\rceil$
plus an absolute constant; affine outputs need no binaries. It also gives
a polynomial rational construction for dense polynomial data
(\cref{thm:vector-rank}). Selecting a basis from original convex outputs
preserves the nonnegative chord gaps needed for this bound. Positive facet
scalarizations and positive polar spanners extend the comparison to the
stated unconditional error bodies, including a strong-separation-oracle
model (\cref{thm:vector-oracle-rank}), where the compiled overhead is
$17+2\lceil\log_2r\rceil$. The constructed optimization model
has an explicit rational inner band. Positive separable polynomial systems
admit allocation comparisons and constructions polynomial in sparse
binary-exponent encoding (\cref{thm:positive-loglog}). Separable vector
maps admit a common curvature-basis comparison, with the rational
construction stated for dense polynomial data
(\cref{thm:separable-vector-rank}).

The limits of these comparisons are part of the results. Fixed-degree
convex product graphs have exact counts $n$ and $\lceil n\log_2 3\rceil$;
a nonconvex scalar polynomial family has two general integers and growing
binary dimension (\cref{thm:exact-box-gap,thm:nonconvex-gap}). A tilted
error body transfers this gap to componentwise convex outputs while its
condition number grows. A root graph has a four-bit formulation but
requires long rational MILP descriptions, whereas a short rational
mixed-integer second-order cone description exists
(\cref{thm:root-encoding,thm:root-conic-separation}). These examples show why
count, rational description, input dimension, and error geometry must be
specified separately.

\subsection{Relationship to prior work}

Table~\ref{tab:result-map} summarizes the closest comparisons. Its last
column describes results proved here; it does not attribute novelty to
all of their ingredients.
\begin{table}[htbp]
\centering\small
\begin{tabular}{@{}>{\raggedright\arraybackslash}p{.25\textwidth}>{\raggedright\arraybackslash}p{.32\textwidth}>{\raggedright\arraybackslash}p{.35\textwidth}@{}}
\toprule
Prior work & Established result & Development in this paper\\
\midrule
\citet{lubin2022} & Minimum MICP rank and midpoint lower bound & Accuracy-dependent contact-volume bounds for all intermediate graph relaxations\\
\citet{beach2022,ggow2020,iqs2018} & Compact quadratic approximations; noncommutative rank and capacity theory & Exact joint quadratic coefficient and uniform finite covariance comparison\\
\citet{rebennack2015,simchowitz2018,codsi2025} & Minimum-breakpoint models, chord estimates, and adaptive or greedy approximation & Two-bit scalar comparison and eleven-bit polynomial rational construction\\
\citet{lyu2026,awerbuch2004} & Common SOS2 encodings and barycentric spanners & Curvature-rank count comparisons with box, facet, and oracle error models\\
\bottomrule
\end{tabular}
\caption{Established methods and the complete graph-approximation comparisons developed here.}
\label{tab:result-map}
\end{table}

\paragraph{Representability and formulation size.}
\citet[Definition 4.3 and Lemma 4.1]{lubin2022} define minimum MICP rank
and show that $w$ pairwise midpoint-incompatible points require at least
$\lceil\log_2w\rceil$ integer variables. We use the same parity mechanism,
then control whole contact classes by covariance or chord geometry.
Taking closures of contacts, rather than of the original lift, makes the
volume arguments valid without measurable witness selections or closed
lifts. On the upper-bound side, \citet[Proposition 1]{vielma2018} embed
finite unions of bounded rational polyhedra using distinct binary codes.
Our elementary version, \cref{lem:disjunction}, also permits real coefficients.
This provides an
encoding once the local graph bands have been chosen; it does not by
itself bound their number or the size of an implicit family.

Mixed-integer extension complexity likewise distinguishes integer count
from linear formulation size. \citet{cevallos2018} derive lower bounds
on integer count when the number of linear inequalities is restricted,
using formulations whose projected mixed-integer hull represents a
combinatorial polytope; \citet{schade2026} strengthen these bounds for
stable set and knapsack. Those size measures count inequalities, rather
than rational encoding bits. Our accuracy-dependent graph model instead
requires the projected set itself to stay in the prescribed error band,
and its minimum integer count allows arbitrary convex lifts without a
size restriction.

\paragraph{Quadratic discretization and matrix-space geometry.}
\citet{beach2022} give quadratic approximations whose binary, row, and
continuous-variable counts grow logarithmically in reciprocal error.
Product envelopes \citep{mccormick1976}, multiparametric disaggregation
\citep{teles2013}, and its refinements \citep{beach2024,beach2024b} supply
related shared discretization tools. Accordingly, logarithmic growth and
binary product linearization are antecedents of our constructions.
\citet{ggow2020} develop operator-scaling and capacity theory, and
\citet{iqs2018} give constructive polynomial-time noncommutative-rank
certificates. We import their algebraic results; the symmetric coordinate
allocation and contact-volume arguments connect them to graph accuracy.

Anisotropic interpolation uses Hessian information to select local element
shapes \citep{cao2007}. Classical determinant maximization treats linear
matrix inequality constraints \citep{vandenberghe1998}. Our covariance
constraints are generally quadratic in $P$, and the asserted approximation
is to minimum integer dimension over arbitrary convex lifts.
Geodesic optimization on positive definite matrices
\citep{sra2013,zhang2016} and convex-body separation/optimization
\citep{gls1981} are established tools. We supply the rational error bounds,
feasibility repairs, and residual certificates needed to use them in the
stated bit-complexity model.

\paragraph{Adaptive approximation and coupled outputs.}
\citet{rebennack2015} formulate minimum-breakpoint continuous
piecewise-linear approximation, including shifted under- and
overestimators. \citet[Lemma 2.1]{simchowitz2018} give the factor-two
chord/midpoint estimate and use a local approximation modulus with endpoint
corrections in convex regression. \citet[Sections 3--4]{codsi2025} reduce
corridor fitting to successive maximal segments. For continuously
differentiable convex corridor boundaries, their maximal-segment routine
uses logarithmically many function and derivative oracle calls in the
reciprocal numerical precision, per segment. Although their general model
permits discontinuous piecewise-linear functions, the convex-corridor
algorithm produces a continuous approximant \citep[Remark 3]{codsi2025}.
These precedents establish
adaptive knots, chord estimates, and greedy segmentation. The additional
claims here compare to arbitrary convex integer lifts and, for the
polynomial classes, control total rational encoding without enumerating
all segments. The compiler uses established continuous encodings of
Boolean circuits \citep{avis2019} and products involving discrete
functions \citep{adams2012}.

\citet[Proposition 1]{lyu2026} explicitly share one SOS2 constraint among
several piecewise-linear functions of the same input. We build on this
possibility, comparing the required common partition with $\pconv$ through
curvature rank and error-body geometry. The determinant selection is a
barycentric-spanner method \citep{awerbuch2004}; approximate spanner
construction \citep{plevrakis2020} and separation/optimization equivalence
\citep{gls1981} are also inherited. For the separation examples, the cap-set
estimate is due to \citet{ellenberg2017}, and difference-of-convex
decomposition is classical \citep{hartman1959}. Our claims concern the
stated graph-contact reductions and formulation separations.

The main novelty assertion is therefore the quadratic characterization
\eqref{eq:intro-ncrank}; the finite and nonlinear contributions are the
specific comparison and encoding theorems above. We do not claim optimality
of every additive constant or a characterization of all smooth and vector
maps. In particular, a universal additive binary/general-integer gap for
one-input convex vectors with growing output dimension and box error
remains open here.

The paper proceeds from contact geometry and quadratic precision laws to
finite quadratic accuracy and rational constructions, then scalar convex
geometry and coupled vector comparisons. Relative error and the
root-power encoding barrier clarify the dependence on the accuracy model
and coefficient representation. Each section states its own hypotheses
and defines the rank or allocation benchmark it uses; these quantities
are not interchangeable.
